\chapter*{Préface}
\addcontentsline{toc}{chapter}{Préface}
\markboth{Préface}{}

Le but de ce document est de permettre aux scientifiques en général, aux
étudiants physiciens et ingénieurs en particulier, d'élaborer des méthodes de
calcul numérique utilisables par l'ordinateur pour résoudre les problèmes de
recherche et de développement.

Il s'agit d'une compilation des notes de cours dispensés par le Professeur
WOAFO Paul de 1994 à 2016. C'est en effet en 1994 que cet enseignement a été
introduit dans les programmes de formation des étudiants de Maîtrise de
Physique au Cameroun (à l'Université de Yaoundé, seule université
camerounaise en ce moment-là). Plus tard, un enseignement de méthodes
numériques a été introduit en année de Licence de Physique. Et les travaux
pratiques étaient effectués au Centre de Calcul (aujourd'hui Centre
Universitaire des Technologies de l'Information).

Année après année, l'enseignement s'est enrichi et a contribué à la formation
des étudiants de Maîtrise, Master et Doctorat camerounais qui pour la plupart
étaient engagés dans des travaux théoriques pouvant avoir une forte composante
numérique. Certains de ces étudiants occupent actuellement des positions
académiques au Cameroun et à l'étranger. Certains d'entre eux se sont inspirés
des éléments de ce cours pour former d'autres étudiants dans les universités
camerounaises. L'importance historique de cet enseignement dans la production
scientifique des physiciens camerounais est donc incommensurable.

Ces notes des cours ont bénéficié des contributions de nombreux
collaborateurs, notamment, Professeur NANA NBENDJO Blaise (Faculté des
Sciences, Université de Yaoundé~I), Dr.~NOUBISSIE Samuel (IUT Fotso-Victor de
Bandjoun, Université de Dschang) et Dr.~FAUTSO Gaétan (École Normale
Supérieure de Bambili, Université de Bamenda).

Il faut noter que l'implémentation des méthodes numériques a beaucoup évolué
ces dernières années avec la mise sur le marché des logiciels comme MATLAB
contenant des codes pour les différentes méthodes numériques. Mais, la plupart
des problèmes de recherche et de développement sont complexes et exigent une
expertise en matière de méthodes numériques (schémas numériques et
algorithmes). De plus, pour utiliser les logiciels existant sur le marché, on
doit être capable de comprendre les différentes étapes afin de savoir comment
et où fournir des informations aux logiciels. Il ne faut pas aussi oublier le
prix d'achat des logiciels.

Il nous a donc paru important de mettre ces notes de cours à la disposition
des étudiants. Cette importance est renforcée par la croissance rapide des
effectifs des étudiants ; ce qui a un effet négatif sur les conditions de
formation.

Le document présente les méthodes numériques de base qui sont dispensées en
année de Licence et de Master de Physique et des études en sciences de
l'ingénieur. Il comporte les méthodes numériques pour :
\begin{itemize}[label=--]
  \item le traitement numérique des mesures expérimentales,
  \item le calcul des intégrales et des dérivées,
  \item la résolution des systèmes d'équations algébriques linéaires et non
        linéaires,
  \item la résolution des équations différentielles et
        intégro-différentielles,
  \item la résolution des équations aux dérivées partielles.
\end{itemize}

Les méthodes de simulation numérique des équations issues des problèmes
d'advection-diffusion, d'écoulement des fluides (équations de Navier-Stokes)
et la dynamique des structures atomiques ne sont pas incluses dans ce
document. Elles pourront éventuellement être présentées en séminaires ou de
manière détaillée aux étudiants de 2\textsuperscript{e} année de Master ou de
Doctorat. Il en est de même des approches Méthodes des Éléments Finis et des
Volumes Finis.

Chaque partie du cours fournit d'une part les développements théoriques des
méthodes numériques, et d'autre part des algorithmes correspondants (prêts à
être traduits dans un langage de programmation tels que le Pascal, le Fortran
ou le C, que les étudiants sont supposés bien maîtriser). Une extension de
l'utilisation de certaines méthodes numériques expliquées dans ce document est
destinée à la programmation des microcontrôleurs qui sont susceptibles de
convertir les instructions du programme en signaux électriques directement
utilisables pour des applications en mécatronique ou dans les systèmes
automatiques ou en électronique pour la génération des signaux.

La synthèse présentée dans ce document a bénéficié des supports didactiques
déjà disponibles pour l'enseignement des méthodes numériques, dont notamment
les livres suivants : S.~Vandergraft, \og Introduction to Numerical
Computations\fg{}, Academic Press (New York, 1978) ; R.L.~Burden,
J.D.~Faires, A.C.~Reynolds, \og Numerical Analysis\fg{}, 2\textsuperscript{nd}
edition, Weber and Schmidt (Boston, 1981) et D.~Euvrard, \og Résolution
numérique des équations aux dérivées partielles\fg{}, Masson (Paris, 1987).

\begin{flushright}
  \textit{Woafo Paul, Professeur}
\end{flushright}
